\section{What a finite certificate says about the field}
\label{sec:field}

This section first specializes a fixed algebraic certificate and
then uses Theorem~\ref{thm:finitecertificate} to treat every finite
nonprice.  Although prices are defined by limits, the strict failure
of a finite configuration has a finite algebraic witness.  This also
shows that the price cone depends only on the characteristic.

\begin{proposition}[Specialization of a finite identity certificate]
\label{prop:finite-specialization}
Fix finitely many tensors with integer coefficients and finitely
many tensor identities, restriction identities or nonvanishing
conditions on local linear maps.  Suppose that choices of complex
coefficients satisfy these conditions.  Then choices in a number
field $K$ satisfy them.  Given such choices, there is a computable
positive integer $D$ with the following property: for every prime
$p\nmid D$, all the same conditions hold after reduction to a
finite extension of $\mathbb F_p$.

One may take any $D$ for which all coefficients and the inverses
of all required nonzero values belong to
$\mathcal O_K[1/D]$.  Thus an explicit witness gives the explicit
exceptional set
\[
 \{p:\ p\text{ is prime and }p\mid D\}.
\]
No discriminant prime need be excluded solely because it ramifies.
\end{proposition}

\begin{proof}
Introduce a variable for every unknown entry of every local map.
Equality of tensors means equality of their finitely many
coordinates, each a polynomial with integer coefficients in these
variables.  Replace each required condition $q\ne0$ by the equation
$qv=1$ with a new variable $v$.  The resulting finite system has
a complex solution, so its ideal is proper over $\mathbb Q$.
It stays proper over $\overline{\mathbb Q}$: if a finite
combination of its generators equalled $1$ over that field,
solving the finite linear system for the coefficients of that
combination would already produce such a combination over
$\mathbb Q$.  The weak Nullstellensatz now gives a solution
over $\overline{\mathbb Q}$.  Its finitely many coordinates
generate a number field $K$.

Every algebraic number becomes an algebraic integer after
multiplication by a nonzero integer.  Choose $D$ divisible by
such integers for all coordinates, including the variables $v$.
This is effective from defining polynomials and expressions for
the coordinates.  All identities then hold in
$\mathcal O_K[1/D]$.  For $p\nmid D$, choose a prime ideal
$\mathfrak p$ of $\mathcal O_K$ above $p$.  Reduction is defined
on $\mathcal O_K[1/D]$, and its residue field is a finite
extension of $\mathbb F_p$.  It preserves the identities and
the equations $qv=1$, so it preserves every required
nonvanishing condition.  This argument uses only reduction
in the ring of integers; ramification changes neither the
existence of $\mathfrak p$ nor the reduction homomorphism.
\end{proof}

The proposition is a finite-presentation argument over
$\operatorname{Spec}\mathbb Z[1/D]$.  The word ``explicit''
refers to the integer extracted from a chosen witness.  An
existence proof for that witness, without its coefficients,
does not itself produce a numerical list of exceptional primes.

\begin{lemma}[Descent of spectral inequalities]
\label{lem:spectral-field-descent}
Let $L/k$ be a finite field extension of degree $d$, and let
$s,t_1,\ldots,t_A$ be tensors on a finite set of $r\geq1$
parties, with coefficients in $k$.  Suppose that over $L$,
$s$ is a sum of restrictions of $t_1,\ldots,t_A$.
Every function $\varphi$ satisfying the spectral axioms of
Definition~\ref{def:price}, equivalently \cite[Definition A.1]{Anthropic2026}, over $k$ then satisfies
\[
 \varphi(s)\leq\sum_{j=1}^A\varphi(t_j).
\]
The extension need not be separable.
\end{lemma}

\begin{proof}
Choose a $k$-basis $e_1,\ldots,e_d$ of $L$ with $e_1=1$,
and a $k$-linear map $\lambda:L\to k$ with $\lambda(1)=1$.
First consider a single restriction
$s=(A_1\otimes\cdots\otimes A_r)t$ over $L$.
For every positive integer $n$, expand the entries of the
local map $A_i^{\otimes n}$ in the chosen basis:
\[
 A_i^{\otimes n}=\sum_{j=1}^d e_j B_{ij},
\]
where each $B_{ij}$ is a matrix over $k$.  Applying $\lambda$
to the coefficients of the tensor identity gives
\[
 s^{\otimes n}=
 \sum_{j_1,\ldots,j_r=1}^d
 \lambda(e_{j_1}\cdots e_{j_r})
 (B_{1j_1}\otimes\cdots\otimes B_{rj_r})t^{\otimes n}.
\]
Each scalar can be absorbed in one local map.  Thus this
is a sum of at most $d^r$ restrictions over $k$.  The
spectral axioms imply
$\varphi(s)^n\leq d^r\varphi(t)^n$.

For the asserted sum identity, take its $n$th tensor
power and expand into the $A^n$ choices of terms.
Apply the same coefficient-extraction argument to each
term, using the same basis of $L$.  Multiplicativity and
the sum inequality give
\[
 \varphi(s)^n\leq
 d^r\sum_{j_1,\ldots,j_n=1}^A
        \prod_{i=1}^n\varphi(t_{j_i})
 =d^r\left(\sum_{j=1}^A\varphi(t_j)\right)^n.
\]
Take $n$th roots and let $n$ tend to infinity.  No trace,
and hence no separability, was used.
\end{proof}

This elementary descent is one aspect of the general
extension-field invariance of asymptotic spectra proved
by Strassen~\cite[Theorem~3.10]{Strassen1988}; see also
\cite[Remark~1.3, p.~7 of arXiv version~3]{cvz-universal}.
It is useful here because wire tensors have integer
coefficients.  Reduction to a finite residue field extension
in Proposition~\ref{prop:finite-specialization} therefore
causes no loss in the spectral inequality over any field
of the same positive characteristic.  Indeed, a field $k$
of characteristic $p$ has an algebraic closure containing
the finite residue field; their compositum is a finite
extension of $k$, to which the lemma applies.

\begin{theorem}[The price cone depends only on characteristic]
\label{thm:price-characteristic}
Assume Theorem~\ref{thm:finitecertificate}.  For any field extension
$K/k$ and any fixed finite party set $V$, the cones of prices over
$K$ and over $k$ are identical.  Consequently two fields of the
same characteristic have the same price cone on $V$.
\end{theorem}

\begin{proof}
Every nonprice over $k$ has a finite cover certificate over $k$.
Extending the local maps to $K$ preserves the certificate and its
strict real length margin.  Thus it remains a nonprice over $K$.

Conversely, let $e$ have a certificate over $K$, with loads
$u,v_1,\ldots,v_j$, tensor power $N$, number of terms $M$, and
spanning tree $T$.  The entries of its finitely many local maps
satisfy a finite system of polynomial equations over $k$.
The associated ideal is proper because it has a solution over
$K$.  A maximal ideal of its finitely generated coordinate
algebra has residue field finite over $k$, by the weak
Nullstellensatz.  Hence the same cover exists over a finite
extension $L/k$, of degree $d$.

Write $r=|V|$, $\tau=\sum_{ab\in T}e_{ab}$ and
\[
 \delta=\langle u,e\rangle-\max_i\langle v_i,e\rangle,
 \qquad \gamma=N\delta-\tau\log_2 M>0.
\]
Raise the $L$-cover to the $s$th tensor power.  The coefficient
extraction in the proof of
Lemma~\ref{lem:spectral-field-descent} converts it to a $k$-cover
with at most $d^rM^s$ terms.  Each is a restriction of a product
of exactly $Ns$ factors among the same $G_{v_i}$.  The new strict
margin is at least
\[
 Ns\delta-\tau\log_2(d^rM^s)
   =s\gamma-\tau r\log_2d,
\]
which is positive for all sufficiently large integers $s$.
This also handles $\tau=0$, $d=1$, and $M=1$.
Theorem~\ref{thm:finitecertificate} now excludes $e$ over $k$.
The nonprice sets, and therefore the price cones, are equal.
Every field contains its prime field, so the final assertion
follows by comparing both fields to that prime field.
\end{proof}

\begin{corollary}[Specialization of a scale-free certificate]
\label{cor:price-specialization}
Suppose a homogeneous inequality in finitely many wire
distances is proved for every price from a fixed finite
collection of tensor identities as in
Proposition~\ref{prop:finite-specialization}, by an argument
valid over every field.  Suppose also that a fixed real
configuration violates this inequality strictly.  Then
that configuration is not a price over any field of
characteristic $p\nmid D$, for the integer $D$ of that
certificate.
\end{corollary}

\begin{proof}
The proposition and the descent lemma give the same
spectral inequalities in characteristic $p\nmid D$.
The assumed field-independent argument converts these
to the same homogeneous inequality for price distances.
The real distances of the configuration and the strict
violation are unchanged; no reduction of real lengths
modulo $p$ is being made.  The configuration therefore
cannot be a price in that characteristic.
\end{proof}

The hypothesis concerning a fixed finite certificate is
substantive.  An infinite family of unrelated identities
can have no common exceptional set: the systems $p x_p=1$,
one for each rational prime $p$, have a simultaneous
solution in characteristic zero, whereas no positive
characteristic satisfies all of them.  Every finite
subsystem does specialize away from finitely many primes.
Thus the proposition alone neither specializes an
arbitrary analytic inequality about the whole spectrum
nor supplies a certificate for a nonprice assertion.
The finite-certificate theorem fills precisely this gap
for a fixed finite configuration.

\begin{theorem}[Every finite nonprice specializes]
\label{thm:nonprice-specializes}
Assume the duality theorem used in
Theorem~\ref{thm:finitecertificate}.  Let $e$ be a
nonnegative, nonzero function on pairs of a fixed
finite party set.  If $e$ is not a price over
$\mathbb C$, then it is not a price over any field
of characteristic $p$ outside a finite set.  A
finite cover witnessing nonprice, expressed over
a number field, gives that set explicitly as
the prime divisors of a common integral
denominator $D$ of its coefficients.
\end{theorem}

\begin{proof}
Theorem~\ref{thm:finitecertificate} gives integral
wire loads $u,v_1,\ldots,v_r$, a positive integer
$N$, a cover with $M$ summands, and a spanning
tree $T$, with strict margin
\[
 N\bigl(\langle u,e\rangle-
           \max_j\langle v_j,e\rangle\bigr)
   >\left(\sum_{ab\in T}e_{ab}\right)\log_2 M.
\]
All prescribed tensors in this cover have
integer coefficients.  Proposition~\ref{prop:finite-specialization}
replaces its local maps by maps over a number
field and then reduces the cover in every
characteristic $p\nmid D$.  If the residue
field is not contained in the chosen
coefficient field $k$, pass to their finite
compositum and use
Lemma~\ref{lem:spectral-field-descent}.
Every normalized spectral point over $k$
therefore obeys the same cover inequality.
The spanning-tree bound and price attainment
in Corollary~\ref{cor:price-attained} turn it
into the weak opposite of the displayed
inequality for every price over $k$.
The strict real margin is unchanged.  Thus
$e$ is not such a price.
\end{proof}

The theorem provides no numerical exceptional
set before a witness has been found.  It also
does not guarantee that two is outside that
set.  Consequently this specialization result
does not decide the square in characteristic
two.

\begin{corollary}[Prices in infinitely many characteristics]
\label{cor:infinitely-many-price-characteristics}
Under the same hypotheses, a fixed finite nonnegative function
that is a price in infinitely many distinct positive
characteristics is a price in characteristic zero.
\end{corollary}

\begin{proof}
A nonprice in characteristic zero is a nonprice over
$\mathbb C$ by Theorem~\ref{thm:price-characteristic}.
Theorem~\ref{thm:nonprice-specializes} then excludes all but
finitely many positive characteristics.  This proves the
contrapositive.
\end{proof}

\begin{proposition}[Bounded certificates lift from infinitely many primes]
\label{prop:bounded-certificates-lift}
Fix a finite party set $V$ and a nonnegative function $e$ on its
pairs.  Fix one discrete finite-cover template in
Theorem~\ref{thm:finitecertificate}, including the loads,
tensor power, terms and dimensions of the local maps, and assume
its real length margin is strictly positive.  If its tensor
equations have a solution in infinitely many distinct positive
characteristics, they have a solution over a number field, and
$e$ is not a price in characteristic zero.

The same conclusion holds if, in infinitely many positive
characteristics, $e$ has such certificates with all their
discrete parameters bounded independently of the characteristic.
\end{proposition}

\begin{proof}
The fixed template is a finite system
$f_1=\cdots=f_l=0$ with $f_i\in\mathbb Z[X_1,\ldots,X_b]$.
If it had no complex solution, the weak Nullstellensatz and
descent of the finite linear equations for its coefficients
would give $1\in(f_1,\ldots,f_l)\mathbb Q[X]$.
Clearing denominators yields a nonzero integer $D$ and
$h_i\in\mathbb Z[X]$ with
\[
 D=\sum_i h_if_i.
\]
In every characteristic $p\nmid D$ this identity rules out a
solution over any field.  This contradicts solvability in
infinitely many distinct positive characteristics.
There is therefore a complex solution, and
Proposition~\ref{prop:finite-specialization} supplies one
over a number field.  The unchanged positive length margin
and Theorem~\ref{thm:finitecertificate} exclude $e$ in
characteristic zero.

With $V$ fixed, bounds on all discrete parameters leave only
finitely many templates: there are finitely many load vectors,
trees, choices of source products, tensor powers, numbers of
summands and map dimensions within those bounds.  One template
therefore occurs in infinitely many of the characteristics.
Its length margin is fixed and positive.  Apply the first part.
\end{proof}

Thus a characteristic-zero price can have nonprice certificates
in infinitely many positive characteristics only if their
discrete complexity is unbounded.  Neither this statement nor
the preceding corollary gives the converse implication from
being a price in characteristic zero to being a price in every
positive characteristic.

\subsection{Uniform neighborhoods and limits across characteristics}

The same finite certificate gives more than specialization of a
single distance array. Its strict margin excludes a neighborhood
whose size is independent of the good characteristic.

\begin{proposition}[Robust specialization of a nonprice]
\label{prop:robust-nonprice-specialization}
Let $e$ be a nonprice over $\C$ on a fixed finite set $V$.
There exist $\epsilon>0$ and a nonzero integer $D$ such that no
nonnegative array $d$ satisfying $\|d-e\|_\infty<\epsilon$ is a
price over any field of characteristic $p\nmid D$.

More precisely, choose a finite certificate from
Theorem~\ref{thm:finitecertificate}, with positive margin
\[
 \gamma=N\bigl(\langle u,e\rangle-
                  \max_j\langle v_j,e\rangle\bigr)
                   -\tau_T(e)\log_2M>0.
\]
One may take
\[
 \epsilon=\frac{\gamma}{2L},\qquad
 L=N\|u\|_1+N\max_j\|v_j\|_1+(|V|-1)\log_2M.
\]
The exceptional integer $D$ is determined after an algebraic
witness for the finite tensor equations has been chosen.
\end{proposition}

\begin{proof}
Specialize the certificate as in
Theorem~\ref{thm:nonprice-specializes}.
For every $p\nmid D$, its same finite cover exists over an algebraic
extension of $\F_p$. The characteristic invariance of prices in
Theorem~\ref{thm:price-characteristic} therefore makes the same
homogeneous cover inequality valid over every field of that
characteristic.

Write its left side minus right side as
\[
 F(d)=N\langle u,d\rangle-N\max_j\langle v_j,d\rangle
                                      -\tau_T(d)\log_2M.
\]
The inequality $F(d)\le0$ holds for each of those prices.
For any arrays $d,d'$,
\[
 |F(d)-F(d')|\le L\|d-d'\|_\infty,
\]
because a maximum of the stated linear forms has Lipschitz
constant at most $\max_j\|v_j\|_1$ and a spanning tree has
$|V|-1$ edges. The positive value $F(e)=\gamma$ implies $L>0$.
For the displayed choice of $\epsilon$, every array in the open
ball has $F(d)>\gamma/2>0$, excluding it from each good-characteristic
price cone.
\end{proof}

\begin{corollary}[Limits of prices in distinct characteristics]
\label{cor:prime-price-limits}
Fix a finite party set $V$. If $p_i$ are distinct primes and
$d_i$ is a price over a field of characteristic $p_i$, and if
$d_i$ converges coordinatewise to $d$, then $d$ is a price over
$\C$. More generally, every compact subset of the complement of
the characteristic-zero price cone, inside the nonnegative arrays,
is disjoint from all positive-characteristic price cones outside
finitely many primes.
\end{corollary}

\begin{proof}
If $d$ were a characteristic-zero nonprice,
Proposition~\ref{prop:robust-nonprice-specialization} would exclude
a fixed neighborhood of $d$ in all but finitely many
characteristics. Eventually $d_i$ belongs to this neighborhood
and $p_i$ is outside that finite set, a contradiction.
For the compact assertion, cover the compact set by the
neighborhoods from the proposition, take a finite subcover,
and take the product of their exceptional integers.
\end{proof}

\begin{corollary}[Uniform transfer of a contraction constant]
\label{cor:contraction-specializes}
Fix $\mu\ge2$ and $1\le s<s'<\sqrt\mu$.
If every characteristic-zero price on
$G=\operatorname{Mat}_{\mu}(\F_2)$ satisfies
$D(\mu)\le sD(1)$, then there is a nonzero integer $D$ such that
\[
 D(\mu)\le s'D(1)
\]
for every price over every coefficient field of characteristic
$p\nmid D$. Thus a fixed contraction parameter transfers with
arbitrarily small positive slack in its constant.
\end{corollary}

\begin{proof}
The case $D(1)=0$ is immediate, since rank-one increments connect
all of $G$ and every price is a pseudometric. Otherwise rescale to
$D(1)=1$. Let
\[
 B=\sqrt{2^{\mu^2}(2^\mu-1)^2}.
\]
There are $2^{\mu^2}(2^\mu-1)^2$ directed rank-one increments.
Their mean squared length is one, so each has length at most $B$.
Every difference of matrices is a sum of at most $\mu$ rank-one
matrices; the triangle inequality bounds the diameter by $\mu B$.

Consider all pseudometric arrays on $G$ with $D(1)=1$, diameter
at most $\mu B$, and $D(\mu)\ge s'$.
This is a compact set of nonnegative arrays. None is a
characteristic-zero price, because $s'>s$.
Corollary~\ref{cor:prime-price-limits} excludes this entire set
from every price cone outside finitely many primes.
The required bound follows after undoing the rescaling.
The exceptional integer is obtained from finitely many strict
certificates in the compact-cover argument; this proof does not
supply a numerical list of excluded primes without those witnesses.
\end{proof}

\subsection{Separating equations and bad primes}

The corresponding statement for a fixed orbit-closure
separation has a particularly simple exceptional integer.

\begin{proposition}[Integral separating equation]
\label{prop:integral-separation}
Let $W$ be a finite free $\mathbb Z$-module with an
algebraic action of $G=\operatorname{GL}_{q,\mathbb Z}$,
and let $x,y\in W$.  If
$y\notin\overline{G_{\overline{\mathbb Q}}x}$,
there is $f\in\mathbb Z[W]$ whose pullback under the
orbit map of $x$ is zero and for which $f(y)\ne0$.
For every prime $p\nmid f(y)$,
\[
 y_{\overline{\mathbb F}_p}\notin
 \overline{G_{\overline{\mathbb F}_p}
 x_{\overline{\mathbb F}_p}}.
\]
\end{proposition}

\begin{proof}
The integral orbit map gives a homomorphism
\[
 \rho_x:\mathbb Z[W]\longrightarrow
 \mathbb Z[g_{ij},\det(g)^{-1}].
\]
After extension to $\overline{\mathbb Q}$ its kernel
is the ideal of the orbit closure: a polynomial
vanishes on that orbit exactly when its pullback
vanishes on the reduced group.  Flatness of the
field extension shows that this kernel is obtained
from the kernel over $\mathbb Q$ by scalar extension.
A separating element in the extended kernel is a
finite linear combination of rational elements.
At least one of them is nonzero at $y$.  Clear
its rational denominators to obtain $f\in\mathbb Z[W]$.
The target ring of $\rho_x$ is torsion-free over
$\mathbb Z$, so vanishing of $\rho_x(f)$ over
$\mathbb Q$ implies its integral vanishing.
Reduce that identity modulo $p$.  The reduced
polynomial vanishes on the reduced orbit and its
closure, but has nonzero value at $y$ whenever
$p\nmid f(y)$.  It is the required separator.
\end{proof}

This proof also makes the nonuniformity precise: for a
growing family of permanent and determinant orbit
closures, the polynomial $f$ and the integer $f(y)$
can depend on both matrix orders.  There is no
single exceptional set for the whole family without
an additional uniform construction.

The converse fails for general orbit separations.
Let $\mathbb G_m$ act on $\mathbb A^1_{\mathbb Z}$
by scalar multiplication, and take $x=2$ and $y=1$.
In characteristic zero the orbit of $x$ is the
nonzero locus and its closure is the whole line,
so $y$ belongs to it.  In characteristic two $x=0$,
whose orbit closure is $\{0\}$, and $y=1$ is
outside.  This example answers the converse
question for specialization of orbit separation;
it does not assert any characteristic-dependent
behavior of the square-price problem.

Finally, the group of labels
$\operatorname{Mat}_{\mu}(\mathbb F_2)$ in the
contracting hypothesis is a finite set with a
specified addition law.  It is not the field of
tensor coefficients.  The same index set occurs
in that hypothesis over $\mathbb C$, over
$\mathbb F_2$, and over every other field.
Its appearance does not by itself decide which
tensor identities hold in characteristic two.
